\question{25}{Considere $n$ elementos, cada uno con un peso y un valor, y un entero $V$.  

Diseñe un algoritmo de programación dinámica que, dados los pesos, los valores y $V$, determine el peso mínimo de un subconjunto de elementos cuyo valor total sea al menos $V$. Suponga que tal subconjunto existe y que cada elemento se selecciona completo o se descarta por completo.

E.g., para \inline{weight = [3,1,4,2]}, \inline{value = [4,2,5,3]} y $V = 8$, el peso mínimo es $6$, tomando los últimos dos elementos (peso $4+2=6$, valor $5+3=8 \ge 8$).
}

\begin{enumerate}
  \item (10 décimas) Defina una relación de recurrencia que resuelva el problema. Explique en palabras qué significa la función y qué significa cada variable. Sea sucinto.
  \item (5 décimas) Dibuje o describa el grafo de dependencias de la relación de recurrencia. ¿Cuánto espacio requiere como mínimo una solución de programación dinámica para este problema? ¿Se puede optimizar?
  \item (10 décimas) Implemente una solución de programación dinámica, ya sea bottom-up (tabulación) o top-down (memoización), en Python, Java, C, C++, JavaScript o TypeScript.
\end{enumerate}

\bigskip
